%-------------------------
% CV Target: Rendezvous Software Developer Full-Stack
% Author: Hafidz Mulia
% Bahasa: Indonesia
% Compile: pdflatex -interaction=nonstopmode "cv.tex"
%-------------------------
\documentclass[a4paper,10pt]{article}

%----------- PACKAGES -----------
\usepackage[utf8]{inputenc}
\usepackage[T1]{fontenc}
\usepackage{lmodern}
\usepackage[top=0.44in, bottom=0.44in, left=0.6in, right=0.6in]{geometry}
\usepackage{titlesec}
\usepackage{enumitem}
\usepackage[hidelinks]{hyperref}
\usepackage{xcolor}
\usepackage{fancyhdr}
\usepackage{tabularx}

%----------- COLORS -----------
\definecolor{accent}{HTML}{0F4C81}

%----------- PAGE STYLE -----------
\pagestyle{fancy}
\fancyhf{}
\renewcommand{\headrulewidth}{0pt}
\renewcommand{\footrulewidth}{0pt}
\fancyfoot[C]{\footnotesize\thepage}

%----------- SECTION FORMATTING -----------
\titleformat{\section}
  {\large\bfseries\color{accent}}
  {}
  {0em}
  {}
  [\titlerule]
\titlespacing*{\section}{0pt}{6pt}{3pt}

%----------- CUSTOM COMMANDS -----------
\newcommand{\cvname}[1]{%
  {\centering\LARGE\bfseries #1\par\vspace{2pt}}%
}
\newcommand{\cvtitle}[1]{%
  {\centering\large\color{accent} #1\par\vspace{3pt}}%
}
\newcommand{\cvcontact}[1]{%
  {\centering\small #1\par\vspace{5pt}}%
}
\newcommand{\cventry}[4]{%
  \noindent\textbf{#1} \hfill \textit{#3} \\
  \textit{#2}
  \ifx&#4& \else \\ \small #4 \fi
  \vspace{2pt}
}
\newcommand{\cvproject}[5]{%
  \noindent\textbf{#1} \hfill \textit{#2} \\
  \textit{\small #3} \\
  {\small #4}
  \ifx&#5& \else \par\noindent{\small \href{#5}{Repository / Demo}} \fi
  \vspace{2pt}
}

\begin{document}

%==================== HEADER ====================
\cvname{Hafidz Mulia}
\cvtitle{Pengembang Perangkat Lunak | Web Full-Stack}
\cvcontact{%
  \href{tel:+6285158335740}{+62 851 5833 5740} \quad\textbar\quad
  \href{mailto:hafidzmuliia@gmail.com}{hafidzmuliia@gmail.com} \quad\textbar\quad
  Surabaya, Indonesia \\
  \href{https://hafmul.site}{hafmul.site} \quad\textbar\quad
  \href{https://github.com/hafidzmulia-its}{github.com/hafidzmulia-its} \quad\textbar\quad
  \href{https://www.linkedin.com/in/hafidz-mulia/}{linkedin.com/in/hafidz-mulia}
}

%==================== RINGKASAN ====================
\section{Ringkasan Profesional}
Pengembang perangkat lunak full-stack dan lulusan Matematika ITS dengan pengalaman membangun aplikasi untuk kebutuhan bisnis, pendidikan, pelaporan, dan kesehatan. Terbiasa mengerjakan frontend, backend, API, dan basis data relasional menggunakan JavaScript/TypeScript, React/Next.js, serta Laravel/PHP. Saat ini mengembangkan aplikasi untuk kasus bisnis di PT Semen Indonesia Logistik menggunakan Docker, Linux, dan Laravel Filament.

%==================== PENGALAMAN ====================
\section{Pengalaman Profesional}
\cventry
  {Pengembang Aplikasi Web}
  {PT Semen Indonesia Logistik (SILOG) -- Departemen TIK, Unit Pengembangan Aplikasi}
  {Sep 2026 -- Sekarang | Penuh waktu}
  {Mengembangkan sistem aplikasi untuk kebutuhan bisnis menggunakan Docker, Linux, dan Laravel Filament.}

\cventry
  {Pengembang Perangkat Lunak Freelance}
  {NechCode -- Layanan Solusi Digital}
  {Feb 2026 -- Sekarang}
  {Menerjemahkan alur kerja klien menjadi aplikasi web, dari penggalian kebutuhan dan perancangan data hingga pengembangan API, frontend, dan penyerahan aplikasi.}

\cventry
  {Pengajar Pemrograman}
  {Timedoor Academy}
  {Jan 2025 -- Sekarang}
  {Mengajarkan dasar pemrograman dan mendampingi siswa saat melakukan debugging serta pemecahan masalah secara terstruktur.}

\cventry
  {Magang Pengembang Full-Stack}
  {SMA IT Al Uswah Surabaya}
  {Jan 2025 -- Feb 2025}
  {Membangun sistem manajemen presensi berbasis Laravel dari penggalian kebutuhan hingga pengujian penerimaan pengguna.}

%==================== PROYEK ====================
\section{Proyek Full-Stack Terpilih}
\cvproject
  {SIDEDIKK -- PWA Skrining Kehamilan}
  {2026}
  {Pengembang Full-Stack | Laravel 13, PHP 8.3, MySQL, Tailwind CSS, Alpine.js, PWA}
  {Membangun alur skrining 20 pertanyaan dengan empat klasifikasi risiko, pencatatan assessment, dan dashboard administrasi dengan pelaporan Chart.js.}
  {https://sidedikk.my.id}

\cvproject
  {e-kos -- Platform Pencarian Web GIS}
  {2025}
  {Pengembang Full-Stack | Next.js, React, TypeScript, PostgreSQL/PostGIS, NextAuth, Leaflet}
  {Mengembangkan pencarian spasial, data relasional, autentikasi, dan alur administrasi untuk pencarian rumah kos.}
  {https://e-kos.vercel.app}

\cvproject
  {Portal TPQ -- Sistem Manajemen Pendidikan}
  {2026}
  {Pengembang Full-Stack | Laravel 12, PHP, MySQL, Blade, Tailwind CSS, Alpine.js}
  {Menerapkan alur pengelolaan siswa, guru, dan pelaporan dalam satu dashboard untuk kebutuhan operasional pendidikan.}
  {https://tpqnh.my.id}

%==================== PENDIDIKAN ====================
\section{Pendidikan}
\cventry
  {S1 Matematika}
  {Institut Teknologi Sepuluh Nopember (ITS), Surabaya}
  {2022 -- 2026}
  {IPK: 3.40 / 4.00 | Skripsi: riset performa query PostgreSQL dan TimescaleDB dengan mekanisme partisi memory-aware berbasis PROADAPT pada data deret waktu 50M+ baris.}

%==================== KEAHLIAN ====================
\section{Keahlian Teknis}
\begin{tabularx}{\textwidth}{@{}l X@{}}
\textbf{Bahasa:} & JavaScript, TypeScript, PHP, SQL, Python, HTML/CSS \\[2pt]
\textbf{Frontend:} & React, Next.js, Laravel Blade, Tailwind CSS, Alpine.js, Vue.js \\[2pt]
\textbf{Backend \& Data:} & Laravel, REST API, MySQL, PostgreSQL, PostGIS, MongoDB, SQL Server \\[2pt]
\textbf{Alur Kerja:} & Git/GitHub, Docker, Linux, Laravel Filament, PHPUnit, Laravel Pint, Vite, debugging \\[2pt]
\textbf{Kekuatan:} & Desain skema relasional, integrasi API, kode yang rapi dan mudah dirawat, penggunaan alat bantu AI secara bertanggung jawab, komunikasi teknis \\
\end{tabularx}

\end{document}
